\section{Operational Limitations \& Boundary Conditions}
\label{sec:limitations}

While FedGuard provides a robust control plane for enterprise federated learning, a rigorous scientific evaluation requires detailing its operational boundary conditions, cryptographic trade-offs, and deployment constraints.

\subsection{Scalability Bottlenecks in Large-Scale Client Topologies}

The threshold Secure Aggregation protocol implemented in FedGuard relies on pairwise X25519 Elliptic Curve Diffie-Hellman (ECDH) key agreement among participating clients. Consequently, round setup incurs a quadratic communication and key generation complexity of $\mathcal{O}(N^2)$ with respect to client count $N$. While highly efficient for cross-silo federated learning across multi-institutional cohorts ($N \le 100$), scaling to cross-device federated learning with tens of thousands of mobile or edge clients ($N \ge 10,000$) introduces severe key exchange and secret share distribution overhead. Future extensions must integrate sub-linear or tree-structured aggregation primitives (e.g., LightSecAgg or FastSecAgg) to reduce key agreement complexity to $\mathcal{O}(N \log N)$.

\subsection{Hardware Root-of-Trust and Attestation Prerequisites}

FedGuard's host integrity verification relies on cryptographic hardware attestation using TPM 2.0 or Trusted Execution Environments (TEEs, e.g., Intel SGX, AMD SEV). The security framework assumes that endorsement keys and Platform Configuration Registers (PCRs) are fused into physical hardware. However, legacy edge devices, virtualized containers, or heterogeneous IoT nodes lacking dedicated TPM chips cannot generate hardware-rooted attestation quotes. In such deployment environments, FedGuard must fall back to software-based X.509 certificate authentication, diminishing resistance against ring-0 operating system compromise and kernel-level tampering.

\subsection{Cryptographic Computational Overhead on Massive Deep Learning Models}

The multi-party cryptographic primitives in FedGuard—specifically tensor fixed-point quantization ($\mathbb{R}^d \to \mathbb{Z}_R^d$), modular addition modulo $R = 2^{32}-1$, and Shamir secret share Lagrange polynomial interpolation—are currently executed on host CPUs. For standard vision architectures (e.g., ResNet-18 with $11.17 \times 10^6$ parameters), cryptographic operations introduce under 1.5 seconds of overhead per round. However, for ultra-large deep learning architectures (e.g., Large Language Models exceeding $10^9$ parameters), CPU-bound modular arithmetic becomes a computational bottleneck, necessitating custom GPU-accelerated CUDA kernels for parallel tensor masking.

\subsection{PKI Certificate Management and Revocation Propagation Latency}

Identity authentication and transport security in FedGuard enforce Mutual TLS (mTLS 1.3) backed by an enterprise Certificate Authority (CA). When a client node is evicted due to policy non-compliance or consent revocation, its X.509 digital certificate is revoked. However, propagating Certificate Revocation Lists (CRLs) or querying Online Certificate Status Protocol (OCSP) responders across distributed, multi-cloud enterprise networks introduces non-zero synchronization latency. During this revocation window, a compromised client with active TLS session tickets might attempt unauthorized control plane queries before revocation entries propagate fully.

\subsection{Privacy Masking vs. Byzantine Robustness Operational Trade-off}

FedGuard highlights an inherent cryptographic trade-off between client data privacy and Byzantine adversary defense. In \texttt{SECURE\_AGGREGATION} mode, client update tensors are cryptographically masked, preventing the aggregator from inspecting individual gradients. Consequently, coordinate-wise robust aggregation algorithms (e.g., Trimmed Mean, Coordinate Median, Krum) that rely on evaluating gradient norm distances cannot operate directly over masked tensors. Conversely, \texttt{BYZANTINE\_RESILIENT} mode inspects unmasked updates to reject poisoned gradients but exposes individual updates to a curious server. While \texttt{HYBRID\_AUDITED} mode mitigates this by negotiating session posture, simultaneously achieving cryptographic non-interactively masked Secure Aggregation and coordinate-wise distance filtering remains an open research challenge.

\subsection{Statutory Compliance Boundaries of Machine Unlearning}

To satisfy legal mandates under GDPR Article 17 (Right to Erasure), FedGuard incorporates an automated machine unlearning protocol (\texttt{FedEraser} checkpoint un-rolling). While empirical evaluations demonstrate rapid accuracy recovery post-revocation, current machine unlearning literature lacks a universally accepted mathematical proof guaranteeing absolute zero information residual across all non-linear deep neural network parameters. Furthermore, regulatory definitions of "data erasure" within trained statistical models remain subject to evolving legal jurisprudence across international jurisdictions.

\subsection{Production Orchestrator Deployment Challenges}

Finally, deploying FedGuard requires integrating control plane hook bindings into the underlying FL framework. While Conclave provides a reference implementation for the Flower framework, integrating FedGuard into custom, non-Python, or legacy proprietary distributed learning orchestrators necessitates writing dedicated gRPC sidecar proxies. Maintaining zero-copy tensor serialization across heterogeneous language runtimes remains a practical engineering challenge.
